\section{Conclusion}

% In this work, we present {\it Inference Scaled GraphRAG}, a novel approach for enhancing multi-hop question answering over knowledge graphs by leveraging inference-time scaling. Our method integrates sequential and parallel scaling strategies into the GraphRAG framework, enabling LLMs to interleave structured reasoning and interaction steps for more effective graph traversal.
% In this work, we introduce \textit{Inference Scaled GraphRAG}, a framework for improving multi-hop question answering over knowledge graphs through inference-time scaling. By combining sequential and parallel compute strategies, our method enables large language models to reason more effectively over graph-structured knowledge.

% We demonstrate that allocating additional compute at inference time, through deeper reasoning chains and majority voting, leads to consistent gains in answer accuracy. Our approach is model agnostic, and provides controllable inference behavior.

% Future work may explore reinforcement learning \citep{gulcehre2023reinforcedselftrainingrestlanguage, zhang2024rest} to optimize traversal policies, explore integration with agent-based systems, and extend our method to other structured reasoning domains.



We present \textit{Inference Scaled GraphRAG}, a framework for enhancing multi-hop reasoning over knowledge graphs by applying inference-time scaling. By combining \textit{sequential scaling} (multi-step reasoning) and \textit{parallel scaling} (majority voting), our method enables LLMs to more effectively traverse structured contexts.

Empirical results on the GRBench benchmark demonstrate a \textbf{64.7\%} improvement over traditional GraphRAG and a \textbf{30.3\%} improvement over previous graph traversal methods. Furthermore, our method improves performance on difficult multi-hop questions from \textbf{15.26\%} to \textbf{31.44\%} more than doubling performance. 

Our approach validates that inference time scaling is a practical and general solution for structured graph reasoning with LLMs. 

\paragraph{Limitations and future work}
Currently, our majority voting mechanism selects actions based solely on frequency, without evaluating answer correctness. This means incorrect but frequent trajectories may dominate. Future work may explore the use of external verifiers to assess both final answers as well as intermediate results. 

Our framework relies on in-context examples to teach the LLM how to explore the graph. Future work may use reinforcement learning to directly optimize the traversal policy \citep{gulcehre2023reinforcedselftrainingrestlanguage, zhang2024rest}. This would allow the model to receive feedback on successful reasoning chains, improve long-horizon planning, and adapt its exploration strategy based on task rewards leading to more robust and generalizable graph traversal behaviors.

Together, these results position inference-time scaling as a foundational capability for enabling scalable and robust graph reasoning with LLMs. 